\documentclass[11pt]{article}
\usepackage[T1]{fontenc}
\usepackage{lmodern,microtype,amsmath,amssymb,booktabs,tabularx,array,xcolor}
\usepackage[letterpaper,top=.82in,bottom=.75in,left=.93in,right=.93in]{geometry}
\usepackage[colorlinks=true,urlcolor=teal]{hyperref}
\definecolor{pine}{HTML}{13575A}
\definecolor{ink}{HTML}{182F41}
\definecolor{graytext}{HTML}{425567}
\pagestyle{empty}
\setlength{\parindent}{0pt}
\setlength{\parskip}{5pt}
\newcolumntype{Y}{>{\raggedright\arraybackslash}X}
\begin{document}
{\sffamily\color{pine}\small\bfseries RESEARCH NOTES\hfill 27 NOVEMBER 2026}\par
\vspace{9pt}
{\sffamily\bfseries\color{ink}\fontsize{26}{29}\selectfont Positive averages and prescribed Fourier zeros}\par
{\large Sharp support, square-frequency gaps, and Dirichlet localization}\par
\vspace{6pt}
A positive average supported on a short interval cannot cancel an arbitrary list of oscillations. The first paper determines the best universal constant in Stefan Steinerberger's support-radius question. If $f\ge0$ has positive integral, has support in an arc of radius $r$ on $\mathbb R/\mathbb Z$, and has zero Fourier coefficients at the positive integers in $A$, then
\[
                    r\ \ge\ \frac{1}{\pi\sqrt2}
                              \left(\sum_{a\in A}a^{-2}\right)^{1/2}.
\]
Equal masses at the two ends of an interval show where the constant comes from. Their Fourier zeros are the odd multiples of one frequency. The proof shows that this example bounds every other list of positive zeros simultaneously for a whole class of decreasing convex functions.

\vspace{4pt}
{\large\sffamily\bfseries\color{pine} The six papers}\par
\renewcommand{\arraystretch}{1.22}
{\small
\begin{tabularx}{\linewidth}{@{}p{.32\linewidth}Y@{}}
\toprule
\textbf{Paper} & \textbf{Conclusion}\\\midrule
Support radius & The exact finite-cardinality inequality, its equality cases, and stability for near-extremizers.\\
Square frequencies & The maximal zero-free interval has order $\min(K+1,N)/N^2$ for the squares from $N^2$ through $(N+K)^2$.\\
All Dirichlet zeros & The exact supremal critical tail coefficient is $\pi d/Q_*$ for a conjugation-closed family with the critical-line property. Endpoint laws obey a degree bound.\\
Three frequencies & A finite rule for arithmetic triples and a certified transcendental gap of length $0.618557753484020625\ldots$ for frequencies $1,4,7$.\\
Finite rings & Squaring on the radical controls the loss of progression-free density in characteristic two when $\mathfrak m^3=0$.\\
Graph towers & The polynomial law for spanning-tree valuations starts within a logarithmic number of levels in the voltage height in every fixed rank.\\
\bottomrule
\end{tabularx}}

\vspace{7pt}
The papers extend specific results in the \href{https://github.com/VladimirReshetnikov/ProveIt/tree/ff96fd44ab160e65fb8a9d861243828d07c156cf/Analysis/GeneralizedRiemannHypothesis}{ProveIt research article}. The third paper uses its Dirichlet critical-line input to identify the prescribed divisor with real Fourier frequencies. The other five arguments are unconditional. These are research proofs. The accompanying programs check the indicated finite calculations; the new theorems have not been formalized in a proof assistant. Each paper gives its own references, assumptions, and remaining questions.
\end{document}
